
Based on both quantitative evaluation results and manual inspection of trajectories, we identify several patterns that are associated with model failures to operate reliably in the environment in our experiments.\footnote{The analysis section presents observational patterns and correlations; causal claims about why models fail require controlled ablations which we leave for future work.}

We analyze these failure modes from four complementary perspectives.

\subsection{Non-scalable Decision-Making Capability}

\begin{table}[t]
\resizebox{\columnwidth}{!}{
\begin{tabular}{lccccccc}
\toprule
\multirow{2}{*}{\textbf{Model}} &
\multirow{2}{*}{\textbf{Context}} &
\multicolumn{2}{c}{\textbf{Easy}} &
\multicolumn{2}{c}{\textbf{Middle}} &
\multicolumn{2}{c}{\textbf{Hard}} \\
& & SKU $\uparrow$ & Cat $\uparrow$ & SKU $\uparrow$ & Cat $\uparrow$ & SKU $\uparrow$ & Cat $\uparrow$ \\
\midrule
DeepSeek-V3.2 (Exp.)
& 128K
& 4.75 & 3.31
& \underline{6.83} & \underline{5.34}
& 6.64 & \underline{5.45} \\

Gemini-3 (Fast)
& \textbf{1M}
& \textbf{8.72} & \textbf{4.34}
& \textbf{13.60} & \textbf{12.43}
& \textbf{21.57} & \textbf{13.56} \\

GLM-4.6
& 200K
& 4.82 & 2.98
& 3.23 & 2.46
& \underline{7.08} & 5.22 \\

OpenAI-5.1 Mini
& 400K
& 3.66 & 1.94
& 4.10 & 4.10
& 6.67 & 4.79 \\

Grok-4.1 Fast
& 200K
& \underline{5.78} & \underline{3.68}
& 4.87 & 3.08
& 5.40 & 3.20 \\

Kimi-K2 (Thinking)
& 256K
& 5.06 & 3.09
& 5.17 & 4.77
& 6.91 & 4.75 \\

Qwen-235B (Thinking)
& 256K
& 4.59 & 3.67
& 3.17 & 2.36
& 5.11 & 4.51 \\

\midrule
\midrule
Heuristic Strategy
& --
& 9.03 & 4.87
& 35.34 & 18.61
& 34.82 & 18.52 \\
\bottomrule
\end{tabular}
}
\caption{SKU and Category counts sold by different models across difficulties.
Context denotes the maximum supported context length of each model.
\textbf{Bolded} indicates the best model; \underline{underlined} indicates the second best (Human excluded).}
\label{tab:sold_sku_cat_context}
\vspace{-3mm}
\end{table}


As shown in Appendix~\ref{appendix:total_results}, models achieve reasonable performance in the Easy setting but exhibit consistent degradation as the environment scales in complexity. To better understand this phenomenon, we analyze both the average number of stock keeping units (SKUs) and product categories sold per day, as well as the number of SKUs and categories explicitly considered during the \emph{Evolving Strategy} phase (Tables~\ref{tab:sold_sku_cat_context} and Appendix ~\ref{appendix:category_analysis}).

Across most models, decision-making capability does not scale proportionally with the size of the environment. Instead, performance remains relatively flat despite a substantial increase in the number of available options. Notably, Gemini-3 (Fast), which supports the largest maximum context window, performs better than other models in this respect in our experiments. This pattern is consistent with the hypothesis that larger context capacity helps retain salient information even when the effective interaction context is constrained, though further investigation is needed to establish causality.

Nevertheless, even the strongest models fail to cover the full decision space in our experiments. This observation suggests that current systems may face challenges in expanding their effective decision-making scope as the environment grows, which is associated with systematic performance degradation in larger and more complex settings.

\subsection{Incomplete Decision-Making Due to Limited Information Coverage}

\begin{table}[t]
\resizebox{\columnwidth}{!}{
\begin{tabular}{lcccccccc}
\toprule
\textbf{Model} &
Supplier &
Inventory &
Return &
Rating &
Price &
Review &
Sales &
History \\
\midrule
DeepSeek-V3.2 (Exp.) & 41.4 & 100.0 & 26.7 & 75.4 & 6.6 & 0.6 & 89.2 & 0.0 \\
Gemini-3 (Fast)     & 69.6 & 100.0 & 41.0 & 82.3 & 67.7 & 6.3 & 87.6 & 2.9 \\
GLM-4.6             & 93.9 & 98.3  & 79.8 & 77.8 & 84.1 & 5.1 & 96.3 & 0.0 \\
OpenAI-5.1 Mini     & 58.8 & 99.0  & 5.1  & 78.6 & 3.9  & 10.6 & 84.7 & 0.3 \\
Grok-4.1 Fast       & 89.9 & 99.8  & 84.0 & 94.0 & 88.6 & 36.0 & 96.4 & 0.3 \\
Kimi-K2 (Thinking)  & 57.2 & 90.4  & 25.9 & 51.4 & 36.2 & 11.0 & 64.0 & 0.5 \\
Qwen-235B (Thinking)           & 76.8 & 78.3  & 31.5 & 58.2 & 19.8 & 23.6 & 74.6 & 0.0 \\
\midrule
Average             & 69.7 & 95.1  & 42.0 & 74.0 & 43.8 & 13.3 & 84.7 & 1.0 \\
\bottomrule
\end{tabular}
}
\caption{
Percentage of days on which each model queries specific information sources when making decisions for SKUs included in the final daily strategy.
Higher values indicate greater reliance on the corresponding data source.
}

\label{tab:tool_usage_easy}
\vspace{-3mm}
\end{table}


We analyze the information sources that models attend to when making operational decisions by examining the data queried for the set of SKUs included in each day’s final strategy. This analysis reveals a clear concentration of attention on a limited subset of signals.

Specifically, most models primarily rely on supplier prices, inventory levels, SKU ratings, and historical sales records when deciding how to operate selected SKUs in Table~\ref{tab:tool_usage_easy}. In contrast, several other critical signals—such as recent customer reviews, return rates, and current selling prices—are consistently underutilized or entirely ignored.



Further correlation analysis shows a positive relationship between the frequency of SKU reviews queries and average daily sales performance in Appendix~\ref{appendix:tool_use_analysis}. This observation aligns with the underlying environment dynamics and is consistent with incomplete information coverage being a factor associated with lower decision quality in our experiments. Overall, these results suggest that models often fail to perform sufficiently comprehensive information gathering, which is correlated with suboptimal operational decisions.
% \begin{table}[t]
% \resizebox{\columnwidth}{!}{
% \begin{tabular}{lccccccc}
% \toprule
% \multirow{2}{*}{\textbf{Model}} &
% \multirow{2}{*}{\textbf{Context}} &
% \multicolumn{2}{c}{\textbf{Easy}} &
% \multicolumn{2}{c}{\textbf{Middle}} &
% \multicolumn{2}{c}{\textbf{Hard}} \\
% & & SKU & Cat & SKU & Cat & SKU & Cat \\
% \midrule
% DeepSeek-V3.2 (Exp.)
% & 128K
% & 7.80 & \underline{4.07}
% & 6.43 & 4.25
% & 4.71 & 3.65 \\

% Gemini-3 (Fast)
% & \textbf{1M}
% & \textbf{8.50} & 3.95
% & \textbf{11.89} & \textbf{11.18}
% & \textbf{13.32} & \textbf{8.09} \\

% GLM-4.6
% & 200K
% & 6.61 & 3.55
% & 6.01 & 3.48
% & \underline{8.47} & \underline{6.05} \\

% OpenAI-5.1 Mini
% & 400K
% & 6.80 & 2.64
% & \underline{6.82} & \underline{6.81}
% & 7.87 & 5.19 \\

% Grok-4.1 Fast
% & 200K
% & \underline{7.88} & \textbf{4.35}
% & 6.51 & 3.53
% & 7.69 & 3.90 \\

% Kimi-K2 (Thinking)
% & 256K
% & 5.92 & 2.94
% & 5.53 & 4.09
% & 6.03 & 3.22 \\

% Qwen-235B
% & 256K
% & 4.78 & 3.66
% & 6.49 & 4.05
% & 5.77 & 4.23 \\

% \midrule
% \midrule
% Heuristic Strategy
% & --
% & - & -
% & - & -
% & - & - \\
% \bottomrule
% \end{tabular}
% }
% \caption{SKU and Category counts observed during the strategy stage across difficulties.
% Context denotes the maximum supported context length of each model.
% \textbf{Bolded} indicates the best model; \underline{underlined} indicates the second best (Heuristic Strategy excluded).}
% \label{tab:strategy_sku_cat_context}
% \vspace{-3mm}
% \end{table}



\begin{table}[t]
\centering
\small
\resizebox{\columnwidth}{!}{%
\begin{tabular}{lcccccc}
\toprule
 & \multicolumn{3}{c}{\textbf{Macro Strategy}} 
 & \multicolumn{3}{c}{\textbf{Execution Strategy}} \\
\cmidrule(lr){2-4} \cmidrule(lr){5-7}
\textbf{Model} 
& Std\_diff~(↓) & MAC~(↓) & TV~(↓)
& Std\_diff~(↓) & MAC~(↓) & TV~(↓) \\
\midrule
DeepSeek-V3.2(Exp.)  
& 0.130 & 0.090 & 5.00 
& 0.289 & 0.144 & 7.93 \\
Gemini-3 (Fast)    
& 0.136 & 0.085 & 3.82 
& 0.293 & 0.225 & 10.18 \\
GLM-4.6          
& 0.131 & 0.096 & 4.52 
& 0.264 & 0.209 & 9.87 \\
OpenAI-5.1 Mini       
& 0.069 & 0.051 & 2.50 
& 0.229 & 0.166 & 8.00 \\
Grok-4.1 Fast    
& 0.079 & 0.045 & 2.71 
& 0.204 & 0.151 & 9.39 \\
Kimi K2(Thinking)
& 0.179 & 0.133 & \textbf{6.95} 
& 0.335 & 0.271 & \textbf{14.08} \\
Qwen-235B (Thinking)        
& \textbf{0.240} & \textbf{0.189} & 6.51 
& \textbf{0.391} & \textbf{0.306} & 10.57 \\
\bottomrule
\end{tabular}
}

\caption{Temporal instability metrics of macro- and execution-level strategy similarity in the Easy environment. All metrics are lower-is-better (↓). Larger values indicate greater temporal instability. Column-wise maximum values are highlighted in bold.}
\label{tab:strategy_instability_easy}
\end{table}
\subsection{Temporal Instability in Execution-Level Decision-Making}

Beyond limitations in decision scalability and information coverage, we observe temporal instability in execution-level decision-making which is associated with long-horizon failures in our experiments. Even under relatively stable environmental conditions, agents frequently revise their strategies across consecutive days, resulting in inconsistent execution trajectories.

\paragraph{Measuring Strategy Similarity.}
We quantify temporal instability by measuring the similarity between strategies on adjacent days at both macro and execution levels. Macro strategy similarity is assessed using an LLM-based prompt that evaluates semantic consistency between consecutive high-level plans. Execution strategy similarity is computed via set-based Jaccard similarity over key fields, including \texttt{focus\_skus}, \texttt{sku\_supplier\_mapping}, \texttt{news\_to\_monitor}, and \texttt{sku\_to\_monitor}, with the final execution similarity obtained by averaging across fields.

\paragraph{Instability Metrics.}
To characterize temporal fluctuations, we compute three complementary metrics: the standard deviation of first-order differences (\textit{Std\_diff}) to capture short-term volatility, the mean absolute change (\textit{MAC}) to estimate typical day-to-day variation, and total variation (\textit{TV}) to reflect cumulative long-term instability.

Across all three environment configurations, we observe consistent temporal instability in both macro- and execution-level strategies. For clarity, we present detailed results from the Easy environment in Table~\ref{tab:strategy_instability_easy} and Appendix~\ref{appendix:score_analysis} . Despite its reduced complexity, the Easy setting already exhibits pronounced temporal fluctuations. In particular, macro strategies remain relatively stable over time, whereas execution strategies show substantially larger variability across most models. This effect is especially pronounced for Qwen-235B (Thinking), which also performs poorly across all environments.

Overall, these results suggest that long-horizon failures in our experiments are associated not only with suboptimal strategy formulation, but also with the inability to maintain temporally consistent execution policies over extended horizons. Causal attribution requires controlled ablations which we leave for future work.
\subsection{Hallucinations and Invalid Actions}

Finally, manual inspection reveals recurrent failure patterns that directly break planning correctness and action validity. We distinguish two closely related but practically different issues:

\paragraph{Hallucinations in reasoning}
Models occasionally generate reasoning traces that reference non-existent SKUs or fabricate numerical quantities, and subsequently incorporate these hallucinated elements into multi-step plans (examples are provided in Appendix~\ref{appendix:hallucinations}). As a result, decisions become misaligned with the true environment state, even when the overall planning structure appears internally coherent.

  
\paragraph{Invalid or irrational actions.}
Models sometimes output actions that violate basic constraints or are inconsistent with historical demand, such as negative order quantities or implausible pricing in Appendix~\ref{appendix:invalid_action}. Even though state-modifying operations are restricted to only a small set of tools, these invalid actions occur with non-negligible frequency and can quickly destabilize the system in long-horizon operation.

In realistic operational settings, both hallucinations and invalid actions would be unacceptable and could directly trigger system collapse.

